\scanpage{102}
Les \deleted{\unsure{objets}} faisceaux de cohomologie de $L_\bullet\overset{L}{\otimes}M_\bullet$ sont appelés \uline{\unsure{hypertor}} de $L_\bullet,M_\bullet$, notation $\operatorname{Tor}_i^{\unsure{\underline O|\mathcal O}}(L_\bullet,M_\bullet)$. \deleted{\unsure{Si} $L_\bullet,\unsure{M_\bullet}$} Comme il s'agit ici de foncteurs \inserted{``exacts''} de catégories triangulées, tout triangle exact en l'un des arguments donne lieu à la suite exacte longue habituelle. En particulier, les $\operatorname{Tor}_i(\ ,M_\bullet)$ définissent un foncteur cohomologique en le premier argument, lorsqu'on prend pour celui-ci un Module ordinaire.

La formation de $L_\bullet\overset{\mathbb L}{\otimes}M_\bullet$ commute au changement de base \uline{plat}, \inserted{$f:X\to Y$}, quand celui-ci est \deleted{\unsure{défini convenablement} : [\unsure{mais aussi}}
\inserted{\[
\boxed{\mathbb Lf^*(M_\bullet\overset{\mathbb L}{\otimes}L_\bullet)\simeq\mathbb Lf^*(M_\bullet)\overset{\mathbb L}{\otimes}\mathbb Lf^*(L_\bullet)}
\]}
Pour définir le foncteur $\unsure{f'|f^*}\unsure{-|=}\mathbb Lf^*$ entre $D^-$, il faut noter que la \uline{platitude} de $f$ implique que $f^*$ transforme quasi-isomorphismes en quasi-isomorphismes. D'autre part, $f^*$ transforme plats en plats, d'où la commutation.
\margin{est-ce vrai dans le contexte des topos généraux ?}
--- En particulier, la formation des $L_\bullet\overset{\mathbb L}{\otimes}M_\bullet$ commute au passage aux fibres, ce qui dans le cas qui nous occupe permettra \inserted{\unsure{qqfois}} un calcul pratiquement trivial des $\overset{\mathbb L}{\otimes}$.
\mspage{101}
